\documentclass[11pt]{article}

% --- arXiv-friendly preamble (standard packages only) ---
\usepackage[utf8]{inputenc}
\usepackage[T1]{fontenc}
\usepackage[margin=1in]{geometry}
\usepackage{booktabs}
\usepackage{amsmath,amssymb}
\usepackage{xcolor}
\usepackage{enumitem}
\usepackage[hidelinks]{hyperref}
\usepackage{cleveref} % must load AFTER hyperref

\newcommand{\code}[1]{\texttt{#1}}

\title{\textbf{When Does a Code Graph Help an Agent?}\\
\large Capability-- and Blast-Radius--Bounded Gains for Change-Impact Tasks}

% TODO(author): confirm author list and affiliation before submission.
\author{
  Tapabrata ``Topo'' Pal\\
  \texttt{topo.pal@gmail.com}\\
  % Affiliation: TODO
}
\date{\today}

\begin{document}
\maketitle

\begin{abstract}
A widespread intuition holds that giving an LLM coding agent a resolved code
graph (callers, callees, implementors, dispatch) makes it more effective on
existing code than plain text search (grep/ripgrep + read). In a controlled,
paired study that isolates the context mechanism and scores against independent
compiler-grade oracles---across two languages (Go, Java), three model tiers, and
change-impact tasks spanning 1--108 change-sites---we find the intuition is
\emph{true but conditional}. On small, name-greppable changes (most of our Go
sample and the small Java tasks) the graph ties text: a capable agent reaches the
same recall by grep+read. On large-blast-radius changes ($\approx$100
change-sites) the graph acts as a \emph{capability equalizer}, lifting weaker
models toward frontier-level completeness (mean recall $+0.17$ to $+0.24$ for a
small model, $+0.13$ to $+0.17$ for a mid model) while collapsing the variance
that makes text search silently incomplete. That recall advantage shrinks as the
model strengthens and, at the frontier, \emph{vanishes} on tasks the model can
fully enumerate by text but \emph{persists} on harder ones---where the graph
remains more reliable and cheaper. The discriminator is blast radius (the number
of change-sites reachable only via caller edges), \emph{not} name ambiguity or
language. We give the mechanism, a deterministic type-resolution oracle and a
size-graded task generator, and a scoring pitfall (line-number normalization)
that silently penalizes graph tools; and we argue the field should evaluate
code-context tools on \emph{completeness-and-consistency at fixed cost,
conditioned on task size}.
\end{abstract}

\paragraph{Honesty note.}
An earlier version of this study reported a \emph{negative} result (``code graphs
don't beat text search''). That conclusion was too broad: it was drawn from Go
tasks that were, in hindsight, almost all small and name-greppable. Extending the
study to a larger, polymorphism-heavy Java framework
(\code{jackson-databind})---with a precise type-resolution oracle and a
size-graded task set---surfaced the boundary the Go data had missed. The honest
result is \emph{conditional}: the graph's value is real but bounded to
large-blast-radius change-impact tasks, and its character changes with model
capability. We report exactly where it helps, where it does not, and the
mechanism that predicts which is which.

\section{Introduction}
\label{sec:intro}

LLM coding agents work on existing code through a loop: \textbf{locate} relevant
code (search), \textbf{read} it, \textbf{traverse} structure (who calls / what
implements), and \textbf{stop} when subjectively confident. Two paradigms supply
the context: text search, and a resolved code graph. The graph holds strictly
more accurate structural information, so the natural hypothesis is that it
produces more complete, better-calibrated behavior on change-impact tasks, where
a missed call site is a broken build.

We set out to confirm this in Go and \emph{could not}---graph tied text on
completeness and calibration. The honest reading at the time was a negative
result. But our Go tasks, after outcome-blind auditing, were almost all
\emph{small} (1--22 change-sites) and \emph{name-greppable}. The question the Go
data could not answer was whether the graph earns its keep where text search
structurally struggles. We therefore extended the study to \code{jackson-databind},
a polymorphism- and reflection-heavy Java framework whose interface methods
(\code{JsonDeserializer.deserialize}, \code{JsonSerializer.serialize}, \ldots)
have dozens of implementations and call sites---and built a \emph{size-graded}
task set with a precise type-resolution oracle. That extension found the boundary.

\paragraph{Contributions.}
\begin{enumerate}[leftmargin=*]
\item \textbf{A methodology that isolates the tool across two languages}---identical
tasks, three arms (text / graph / graph-as-verifier), enforced tool allowlists,
scored against \emph{independent} oracles (Go \code{go-ssa-vta}; a new Java Spoon
oracle that resolves every call site by type, not by name).
\item \textbf{The conditional result}: graph $\approx$ text on small/greppable
change-impact tasks; graph $\gg$ text on large-blast-radius tasks, on recall
\emph{and} variance.
\item \textbf{The discriminator}---blast radius, established against a confound we
expected to matter (name ambiguity) and which did not.
\item \textbf{A capability-conditioned reading}: the graph is a \emph{quality}
lever for weak/mid models on large changes and a \emph{cost} lever at the
frontier; one tool, two payoffs.
\item \textbf{A reproducible Java change-impact oracle + size-graded task
generator}, and a scoring fix (line$\rightarrow$method normalization) without
which the graph arm is silently under-credited.
\end{enumerate}

\section{Related work}
\label{sec:related}

% TODO(author): citation details below are best-effort and MUST be verified
% before submission (see paper/README.md). Cross-listed cs.SE / cs.AI.

\paragraph{Agentic coding and repository-level evaluation.}
Benchmarks such as SWE-bench~\cite{swebench} and repository-level completion
suites~\cite{repobench,repocoder} established that resolving real changes requires
\emph{repository} context, not just local context. Agent scaffolds---SWE-agent's
agent--computer interface~\cite{sweagent}, and pipeline approaches such as
Agentless~\cite{agentless}---differ chiefly in \emph{how} they gather that
context. Our study isolates one axis of that design space (text search vs.\ a
resolved code graph) and measures it directly, rather than as part of an
end-to-end agent.

\paragraph{Code context: retrieval vs.\ structure.}
Retrieval-augmented generation~\cite{rag} and its code-specific
variants~\cite{repocoder} supply context by similarity; code-graph approaches
supply it by resolved structure. Graph-based code context for
LLMs~\cite{codexgraph} and graph-structured retrieval more
broadly~\cite{graphrag} argue structure should help. We ask the prior question
empirically---\emph{when} does the structure change the answer rather than just
the path to it---and find the answer is conditioned on task size and model
capability.

\paragraph{Static call-graph construction (our oracle, not the tool under test).}
Sound change-impact ground truth requires resolved dispatch. We use Class
Hierarchy Analysis~\cite{cha} and Variable Type Analysis~\cite{vta} lineage call
graphs (Go \code{go-ssa-vta}) and, for Java, source-level type resolution via
Spoon~\cite{spoon} (cf.\ bytecode frameworks Soot~\cite{soot} and
WALA~\cite{wala}). Crucially the oracle is \emph{independent} of the graph tool
under test, so we never grade the tool against itself.

\section{A variable framework for agent code-context}
\label{sec:framework}

Three cost axes---\textbf{tokens, latency, round-trips, setup}; six quality
axes---\textbf{completeness, precision, calibration, consistency, freshness,
determinism}; one outcome---\textbf{task success}; one
meta-variable---\textbf{cost-of-error}. Our prior belief was that the graph's
value lived in the quality axes uniformly. The result below is that it lives in
\textbf{completeness + consistency, but only past a task-size threshold}, and at
near-zero extra cost---so the right figure of merit is
\emph{quality-at-fixed-cost conditioned on blast radius}.

\section{Study design}
\label{sec:design}

\paragraph{Arms (only the tool description/allowlist differs).}
\textbf{T} text (rg/grep/read); \textbf{G} graph (typed \code{prism} call-graph
CLI for traversal + rg for anchor discovery); \textbf{V} text-primary with
graph-as-verifier. Arm enforcement is by \code{--allowedTools} with a recorded
\code{tool\_trace}; runs where the G arm never touched the graph are flagged
(\code{violation}) and excluded.

\paragraph{Mode A (context quality).}
The agent answers ``list every site that must change to do X,'' declaring
\code{complete} + \code{unresolved}. This isolates the tool from patch-writing
skill. (Mode B / compile-and-test is future work, \cref{sec:threats}.)

\paragraph{Independent oracles---never the graph under test.}
\emph{Go:} the \code{go-ssa-vta} call graph (via \code{grove-eval}) and/or the
merged-PR diff (codegen/mocks filtered). \emph{Java:} a Spoon type-resolution
oracle. Given a target method, it emits the declaration + the full
override/implementation family + every call site whose executable \emph{resolves}
(by type, not by name) into that family. This is the discipline a bare-name
oracle lacks: a call to an unrelated same-named method on a different type is
excluded. On jackson it resolved all $\sim$500 \code{get}/\code{set} call sites
and kept only the 3--22 that are truly jackson's (0 unresolved). Ground truth is
human-spot-checked and, where text reaches it, validated by a text arm hitting
recall 1.0.

\paragraph{A scoring pitfall (Java).}
\code{prism} reports authoritative \code{file:line}, so graph-arm agents naturally
answer sites as \code{File.java:114} rather than \code{File.java:method}. A
name-based scorer matches those at 0, silently penalizing the \emph{graph} arm
(e.g.\ a 148-site correct answer scored 0.0). We build a Spoon
line$\rightarrow$method index and normalize numeric answers to their innermost
enclosing method before scoring. All Java numbers below are post-normalization.
We flag this because it is an easy way to accidentally under-report a structural
tool's performance.

\paragraph{Subjects and size grading.}
Go: 14 tasks, 1--93 sites (localization controls, greppable enumeration controls,
change-impact). Java: \code{jackson-databind} @ 2.18.8, 6 interface-method
signature-change tasks spanning \textbf{8, 22, 38, 58, 104, 108} change-sites,
chosen to trace a size curve and to vary name-ambiguity independently of size.
\textbf{Models:} a small, a mid, and a frontier tier (denoted Haiku / Sonnet /
Opus). \textbf{Trials:} 5 per cell; we report tails (min, variance), not just
means. The full grid is 157 scored runs.

\section{Results}
\label{sec:results}

\subsection{Small / greppable change-impact: graph $\approx$ text}
Across the Go tasks, mean recall(G) $\approx$ recall(T) (within $\sim$0.01--0.07
every task; no consistent direction). Greppable controls reach recall $\approx$1.0
for every arm. The small Java tasks behave identically: \code{jsonnode-get} (8
sites) and \code{settable-set} (22 sites) tie ($\Delta$recall $\le 0.01$) on the
small model---\emph{despite} being the most name-ambiguous targets in the set
(grep over-match $64\times$ and $18\times$). Ambiguity did not help the graph.

\subsection{The size curve}
\label{sec:curve}
\Cref{tab:curve} gives the full size-graded curve (recall, $n=5$,
post-normalization). The small model ran the 4 original tasks; the mid and
frontier models ran all 6.

\begin{table}[ht]
\centering
\small
\begin{tabular}{lrccccc}
\toprule
task & sites & Small T & Small G & Mid T & Mid G & Frontier T \(\to\) G \\
\midrule
jsonnode-get      & 8   & 0.97 & 0.97 & 0.78 & \textbf{0.93} & 0.93 \(\to\) 0.95 \\
settable-set      & 22  & 0.90 & 0.91 & 0.93 & 0.97          & 0.97 \(\to\) 0.99 \\
writeTypePrefix   & 38  & ---  & ---  & 0.97 & 1.00          & 1.00 \(\to\) 1.00 \\
serializeWithType & 58  & ---  & ---  & 0.99 & 1.00          & 0.99 \(\to\) 1.00 \\
deserialize       & 104 & 0.52 & \textbf{0.69} & 0.76 & \textbf{0.93} & 0.92 \(\to\) \textbf{0.96} \\
serialize         & 108 & 0.62 & \textbf{0.85} & 0.86 & \textbf{1.00} & 1.00 \(\to\) 1.00 \\
\bottomrule
\end{tabular}
\caption{Recall across three model tiers and the jackson size curve. The graph's
recall advantage is largest for the weakest model on the largest tasks, shrinks
as the model strengthens, and is $\le 0.04$ everywhere at the frontier.}
\label{tab:curve}
\end{table}

The robust, large signal is at the \textbf{large tasks (104/108): graph $+0.13$ to
$+0.24$, on both the small and mid tiers.} The \textbf{mid-size dispatch tasks
(38/58) tie on recall}---the mid model's text arm maxes out (0.97--0.99) because
their change-sets are small and the call sites are name-matched. And the advantage
is \emph{not} explained by grep ambiguity, which runs opposite to the effect (the
most ambiguous targets are the smallest).

\paragraph{The one small-task exception is mechanistically informative.}
On \code{jsonnode-get} (8 sites) the mid model's \emph{text} arm drops to 0.78,
and the graph recovers it ($+0.15$). The sites text consistently missed are
\code{JsonNode.has}, \code{JsonNode.hasNonNull}, and \code{ArrayNode.\_at}---methods
that \emph{call} \code{get(int)} but whose own names do \emph{not} contain
``get.'' A grep-for-\code{get} agent finds the get-named declarations and call
lines but misses \emph{callers named after something else}; the graph's caller
edges return them directly. So the true predictor is not raw size but
\textbf{the number of change-sites reachable only via caller edges, not by the
target's name}---which grows with blast radius (the large tasks have many such
indirect callers) but can bite even a tiny task. This is the sharpest single
illustration of the mechanism.

\subsection{The gain survives a stronger model, then converts at the frontier}
\label{sec:tiers}
A natural objection: a stronger model brute-forces the large enumeration with text
and closes the gap. Going small$\rightarrow$mid, text \emph{did} improve
(serialize $0.62\rightarrow0.87$)---but the graph improved more
($0.85\rightarrow0.996$) and stayed ahead on \emph{both} large tasks
(\cref{tab:tiers}). At the frontier the advantage is conditioned on task
difficulty: on \code{serialize} the frontier text arm fully enumerates the change
(1.00, 24 turns) and the graph is redundant; but on the harder \code{deserialize},
even frontier text drops sites (mean 0.923, one run at \textbf{0.71}) while the
graph holds 0.962---higher recall, lower variance (sd .05 vs .12), \emph{and}
cheaper (\$4.69 vs \$5.64/run). The recall advantage shrinks monotonically with
capability but \emph{does not vanish on hard tasks}---it converts into a
reliability + cost edge.

\begin{table}[ht]
\centering
\small
\begin{tabular}{lcc}
\toprule
tier & serialize (108) & deserialize (104, \emph{harder}) \\
\midrule
Small    & $.62 \to .85$ (\textbf{+.24}) & $.52 \to .69$ (\textbf{+.17}) \\
Mid      & $.87 \to 1.00$ (\textbf{+.13}) & $.76 \to .94$ (\textbf{+.17}) \\
Frontier & $1.00 \to 1.00$ (+.00)         & $.92 \to \mathbf{.96}$ (\textbf{+.04}) \\
\bottomrule
\end{tabular}
\caption{Recall (text $\to$ graph, $\Delta=$ G$-$T) on the two large tasks across
tiers. The graph is a capability equalizer: large lift for weak models, narrowing
to a hard-task reliability/cost edge at the frontier.}
\label{tab:tiers}
\end{table}

\subsection{The graph removes catastrophic misses (a consistency gain)}
\label{sec:consistency}
The tier data show more than a higher mean: the \emph{variance} differs. Mid-model
text on \code{serialize} swings from 1.00 to 0.73 across nominally identical
runs---it silently misses $\sim$30\% of sites on some trials. The graph is
$0.996 \pm {\sim}0$: \emph{zero catastrophic runs}. On a change task, where one
missed site is a broken build, this reliability is arguably more valuable than the
mean lift. This is the completeness/calibration benefit the single-tier Go study
looked for and could not find; it emerges only once tasks are large enough that
text \emph{can} fail, and it persists to the frontier on hard tasks (frontier text
dipped to 0.71 on \code{deserialize}; the graph stayed at 0.96).

\subsection{Cost: cheaper in the middle, par at the top, overhead tax at the bottom}
\label{sec:cost}
Cost (USD/run) depends on the regime. On the \textbf{large tasks} (mid tier) the
graph buys $+0.13$--$0.17$ recall for $\sim$10\% more (G/T $\approx$ 1.07--1.16):
``better answer, $\sim$same price''---text cannot reach the graph's recall at
\emph{any} reasonable budget. On the \textbf{mid-size tasks} recall ties but the
graph is $\sim$25--37\% \emph{cheaper} (G/T $\approx$ 0.63--0.77), recovering a
token-efficiency result at matched quality. On the \textbf{small tasks} per-call
overhead makes the graph pricier (G/T $\approx$ 1.24--1.28) where text was already
near-complete. \textbf{At the frontier the graph becomes a cost lever}: recall
ties on all six tasks but the graph is cheaper on 4 of 6, including every mid/large
task (G/T $\approx$ 0.64--0.90), undercutting a frontier model grinding 24--39 text
turns.

\subsection{Calibration}
In aggregate over small tasks, over-confidence is tied (both arms assert
\code{complete} at $\sim$0.9 recall on \code{settable-set}: 5/5 each)---the Go
negative holds there. The graph's calibration benefit is specifically the variance
collapse on large tasks (\cref{sec:consistency}), not a uniform over-confidence
reduction.

\section{Why---the mechanism}
\label{sec:mechanism}

A change to a method's signature forces every declaration/override \emph{and}
every call site to change. Two regimes:
\begin{itemize}[leftmargin=*]
\item \textbf{Few sites, statically named (Go, small Java).} A call
\code{x.Foo(...)} and an implementor \code{func (T) Foo(...)} both contain
\code{Foo}; with $\le\!\sim$20 sites a capable agent greps the name, reads to
confirm, and enumerates them all. The graph's structural precision saves reading
effort (a token wash) but rarely changes \emph{what it finds}. Hence the tie.
\item \textbf{Many sites (Java framework interfaces, $\sim$100).} Now grep+read is
a manual graph traversal at scale: the agent must find every implementor and every
polymorphic call, hold them in context, and not lose any across 50--65 turns. It
is good on average but \emph{unreliable}---it drops sites stochastically,
especially callers not named after the target. The resolved graph enumerates the
family and the resolved callers in one authoritative pass, so completeness is both
higher and consistent.
\end{itemize}

This predicts the boundary: the graph earns its keep wherever the change-set is too
large to enumerate reliably by hand---large refactors, framework-wide interface
changes---and ties text on the small, local edits that dominate everyday work.
Name ambiguity is a red herring at these sizes: a capable agent disambiguates
\code{get}/\code{set} by reading, as long as there are few sites to check.

\section{Discussion}
\label{sec:discussion}
\begin{itemize}[leftmargin=*]
\item \textbf{Evaluate code-context tools conditioned on blast radius.} ``Graphs
help'' and ``graphs don't'' are both wrong; the honest figure of merit is
completeness-and-consistency at fixed cost, as a function of \#change-sites.
\item \textbf{The graph's value morphs with model capability.} On large changes it
is a quality/reliability lever for weak and mid models and a cost lever at the
frontier. One tool, two payoffs, selected by how capable the model already is.
\item \textbf{Language matters only through task size.} Java surfaced the effect
not because Java is ``harder to grep'' but because framework interface methods
produce 100-site change-sets that Go application code in our sample rarely did. The
prediction is testable in large Go codebases with wide interfaces.
\end{itemize}

\section{Threats to validity}
\label{sec:threats}
\begin{itemize}[leftmargin=*]
\item \textbf{Two subjects per language; one framework for the large-Java regime.}
The large-task win rests on \code{jackson-databind} interface methods; replication
on another framework (Spring, Guava) and in a large Go codebase with wide
interfaces is the key external-validity gap.
\item \textbf{Mode A only.} We score the \emph{answer}, not an applied patch. Mode
B (compile / fail-to-pass)---where an incomplete change actually breaks the
build---would test whether the reliability gain converts to task success.
\item \textbf{One graph implementation} and a soft tool gate (the G arm can decline
to use the graph; we exclude those runs, but the offer-vs-enforce distinction is
unmodeled).
\item \textbf{Construct / selection.} Java ground truth is type-resolved (no
bare-name pollution); an earlier bare-name Java attempt was discarded after audit.
Go ground-truth pollution and codegen churn were mitigated after discovery; earlier
positive readings on hand-picked tasks did not survive outcome-blind sampling.
\item \textbf{Determinism.} Recall/precision come from the oracle scorers and cost
from per-run accounting; no LLM is in the measurement loop, so the numbers are
reproducible from the released run logs.
\end{itemize}

\section{Conclusion}
\label{sec:conclusion}
The intuition that a resolved code graph makes an agent more complete on
change-impact tasks is \emph{conditionally true, and the condition is capability
$\times$ blast radius}. On small, name-greppable changes (most of Go, everyday
edits) the graph ties text at every model tier. On large-blast-radius changes it
acts as a \emph{capability equalizer}: it lifts weaker/cheaper models toward
frontier-level completeness with far lower variance. That recall gain shrinks as
the model strengthens and, at the frontier, vanishes on tasks the model can fully
enumerate by text yet persists on harder ones---where the graph still delivers
higher recall, removes the tail risk, and even costs less. The earlier blanket
negative was an artifact of a small-task sample; the blanket positive (``graphs
obviously help'') is equally wrong. The defensible, mechanistic claim: \textbf{a
code graph is a completeness-and-reliability tool that pays off in proportion to
how far a change exceeds what the model's text search can reliably enumerate---large
for weak models, narrowing to hard-task reliability at the frontier, and a wash for
small local edits.}

\section*{Artifact availability}
\label{sec:artifacts}
The graph engine (Grove), the agent-facing context CLI (Prism), and the evaluation
harness (arms runner, tasks, the Spoon change-impact oracle, the
line$\rightarrow$method rescorer, and the recall/cost aggregator) are open source.
All analyses are reproducible \emph{without} an LLM from the released run logs via
the deterministic scorers. % TODO(author): insert repository URLs and a frozen
% release tag / DOI (e.g. Zenodo) for the camera-ready.
\begin{itemize}[leftmargin=*]
\item Grove (code-intelligence graph): \url{https://example.com/grove} % TODO URL
\item Prism (agent context CLI): \url{https://example.com/prism} % TODO URL
\item Evaluation harness + Java oracle + tasks: \url{https://example.com/harness} % TODO URL
\end{itemize}

% ---------------------------------------------------------------------------
% References. NOTE: best-effort entries; verify every field before submission.
% ---------------------------------------------------------------------------
\begin{thebibliography}{99}
\small

\bibitem{swebench} C.~E. Jimenez, J.~Yang, A.~Wettig, S.~Yao, K.~Pei, O.~Press,
and K.~Narasimhan. SWE-bench: Can Language Models Resolve Real-World GitHub
Issues? \emph{ICLR}, 2024.

\bibitem{sweagent} J.~Yang, C.~E. Jimenez, A.~Wettig, K.~Lieret, S.~Yao,
K.~Narasimhan, and O.~Press. SWE-agent: Agent-Computer Interfaces Enable Automated
Software Engineering. \emph{NeurIPS}, 2024.

\bibitem{agentless} C.~S. Xia, Y.~Deng, S.~Dunn, and L.~Zhang. Agentless:
Demystifying LLM-based Software Engineering Agents. \emph{arXiv:2407.01489}, 2024.

\bibitem{repobench} T.~Liu, C.~Xu, and J.~McAuley. RepoBench: Benchmarking
Repository-Level Code Auto-Completion Systems. \emph{ICLR}, 2024.

\bibitem{repocoder} F.~Zhang, B.~Chen, Y.~Zhang, J.~Keung, J.~Liu, D.~Zan,
Y.~Mao, J.-G. Lou, and W.~Chen. RepoCoder: Repository-Level Code Completion
Through Iterative Retrieval and Generation. \emph{EMNLP}, 2023.

\bibitem{rag} P.~Lewis et al. Retrieval-Augmented Generation for
Knowledge-Intensive NLP Tasks. \emph{NeurIPS}, 2020.

\bibitem{codexgraph} X.~Liu et al. CodexGraph: Bridging Large Language Models and
Code Repositories via Code Graph Databases. \emph{arXiv:2408.03910}, 2024.

\bibitem{graphrag} D.~Edge et al. From Local to Global: A Graph RAG Approach to
Query-Focused Summarization. \emph{arXiv:2404.16130}, 2024.

\bibitem{cha} J.~Dean, D.~Grove, and C.~Chambers. Optimization of Object-Oriented
Programs Using Static Class Hierarchy Analysis. \emph{ECOOP}, 1995.

\bibitem{vta} V.~Sundaresan, L.~Hendren, C.~Razafimahefa, R.~Vall\'ee-Rai,
P.~Lam, E.~Gagnon, and C.~Godin. Practical Virtual Method Call Resolution for
Java. \emph{OOPSLA}, 2000.

\bibitem{spoon} R.~Pawlak, M.~Monperrus, N.~Petitprez, C.~Noguera, and
L.~Seinturier. Spoon: A Library for Implementing Analyses and Transformations of
Java Source Code. \emph{Software: Practice and Experience}, 46(9):1155--1179, 2016.

\bibitem{soot} R.~Vall\'ee-Rai, P.~Co, E.~Gagnon, L.~Hendren, P.~Lam, and
V.~Sundaresan. Soot---a Java Bytecode Optimization Framework. \emph{CASCON}, 1999.

\bibitem{wala} IBM. T.J. Watson Libraries for Analysis (WALA).
\url{https://github.com/wala/WALA}.

\end{thebibliography}

\end{document}
